\chapter{Governação, Privacidade e Responsabilidade na União Europeia}
\label{ch:governacao}

Uma equipa pode ter bons testes e continuar sem saber se pode usar os dados, quem responde pelo sistema ou como um cliente contesta uma decisão. São perguntas diferentes da taxa de acerto. Chegam frequentemente tarde porque não parecem problemas de programação. Continuam a ser problemas do projeto.

Este capítulo distingue enquadramento legal, orientação voluntária e escolhas de engenharia. A verificação documental termina em \textbf{3 de novembro de 2026}. Um caso concreto depende da finalidade, dos dados, do papel da organização e das regras setoriais e nacionais aplicáveis.

\section{Um inventário antes dos princípios}
Para cada sistema, regista finalidade, utilizadores, dados, fornecedores, ferramentas, decisões afetadas e responsável. Acrescenta onde corre, quem pode alterar o modelo e o que acontece quando falha. Um inventário pequeno pode ser mais útil do que cinquenta páginas a afirmar que a empresa valoriza a ética.

Na Lavandaria, «preparar um rascunho» e «aprovar uma compensação» têm efeitos diferentes. A mudança exige rever dados, autorização, critérios e responsabilidade. Não é apenas acrescentar uma ferramenta.

\section{RGPD: o tratamento precisa de fundamento}
O RGPD estabelece princípios como finalidade, minimização, exatidão, limitação da conservação e segurança. O consentimento é um fundamento possível de licitude, não uma solução universal; categorias especiais e transferências exigem análise adicional. Pseudonimização não equivale a anonimização. Uma avaliação de impacto pode ser obrigatória quando o tratamento é suscetível de implicar elevado risco~\citep{rgpd2016}.

O artigo 22 trata de decisões exclusivamente automatizadas com efeitos jurídicos ou efeitos similarmente significativos, com condições e exceções. Não proíbe toda a automação nem transforma qualquer clique humano em supervisão suficiente. Prevê um percurso efetivo de intervenção e contestação nos casos aplicáveis~\citep{rgpd2016}.

Como prática de engenharia, começa com dados fictícios e usa apenas os campos necessários. Não assumes que uma conta paga impede treino, retenção ou acesso por terceiros: verifica termos e configuração do serviço.

\section{AI Act: finalidade, papéis e calendário}
O Regulamento (UE) 2024/1689 estabelece obrigações que variam por prática, risco e papel, incluindo regras para modelos de finalidade geral~\citep{aiact2024}. Um chatbot não é automaticamente de alto risco por usar um LLM; finalidade e contexto contam.

Segundo a cronologia oficial do Conselho consultada, grande parte das regras começou a aplicar-se em 2 de agosto de 2026. Alterações adotadas em 2026 deslocaram marcos de alto risco para 2 de dezembro de 2027 e 2 de agosto de 2028~\citep{conselho2026}. Confirma alterações, âmbito e disposições transitórias antes de aplicar datas. Transparência e literacia também devem ser analisadas no âmbito relevante.

A consulta automatizada do texto consolidado encontrou uma barreira de acesso; o registo de fontes documenta-a. A cronologia institucional sustenta as datas acima. Uma revisão jurídica do texto em vigor continua necessária antes de publicar orientação legal ou decidir conformidade de um produto.

\section{Uma checklist não é uma certificação}
O NIST AI RMF e o perfil de IA generativa oferecem um quadro voluntário de gestão de riscos~\citep{nist2024}. OWASP ajuda a identificar falhas. Estas referências têm funções próprias; segui-las não prova conformidade legal.

Em engenharia, ajudam a perguntar que danos são plausíveis, que evidência existe, como se controla acesso e que risco residual é aceite. A resposta deve ficar ligada a quem decide e à versão do sistema, não a uma declaração abstrata.

\section{Supervisão que possa alterar o resultado}
O revisor precisa de informação, tempo, competência e poder para rejeitar a proposta. Se recebe cem decisões por minuto e só pode carregar em «aceitar», a presença humana pouco acrescenta. Mede carga, divergências, erros apanhados e tempo de resposta.

Um modelo não é uma pessoa a quem a organização possa atribuir a culpa. O responsável pelo processo tem de conhecer limites, aprovar o uso e manter alternativa quando o sistema fica indisponível.

\begin{promptlabbox}{Uma ficha de governação que cabe numa reunião}
\textbf{Objetivo e contexto:} avaliar passar de rascunho para envio automático na Lavandaria fictícia. O exercício não certifica conformidade jurídica.

\textbf{Instruções:} preenche finalidade, dados indispensáveis, efeitos, responsável, contestação, retenção e incidente. Identifica decisões que precisam de jurídico, proteção de dados ou dono do processo.

\textbf{Resultado esperado:} uma ficha com decisões e questões abertas. Um bloqueio relevante deve impedir a funcionalidade afetada, não desaparecer numa nota.

\textbf{Validação:} outra pessoa deve conseguir dizer quem pode parar o sistema e como um cliente corrige um erro. Confirma interpretações nas fontes oficiais com o responsável competente.

\textbf{Falhas possíveis:} copiar uma política genérica, confundir consentimento com toda a autorização ou chamar supervisão a uma aprovação sem informação.
\end{promptlabbox}

Guarda versão aprovada, testes, limites e quem aceita o risco residual. Define reavaliação para novo fornecedor, ferramenta, dados ou finalidade. A governação acompanha mudanças concretas: não precisa de bloquear pequenas experiências, mas tem de reconhecer quando elas passam a afetar outras pessoas.
